\documentclass[11pt]{article}

\usepackage[T1]{fontenc}
\usepackage[utf8]{inputenc}
\usepackage[margin=0.75in]{geometry}
\usepackage{enumitem}
\usepackage{titlesec}
\usepackage{amsmath}
\usepackage{amssymb}
\usepackage{hyperref}

\hypersetup{
    colorlinks=true,
    linkcolor=black,
    urlcolor=black,
    citecolor=black,
    pdfauthor={Miraj Chamara Samarakkody},
    pdftitle={Curriculum Vitae - Miraj Chamara Samarakkody},
    pdfsubject={Curriculum Vitae (Short)},
    pdfkeywords={mathematics, manifold learning, formal proof verification, differential geometry, Lean 4}
}

\titleformat{\section}{\normalsize\bfseries\uppercase}{}{0pt}{}[\titlerule]
\titlespacing*{\section}{0pt}{8pt}{4pt}

\setlist[itemize]{leftmargin=*,topsep=1pt,itemsep=1pt,parsep=0pt}

\pagestyle{empty}
\setlength{\parindent}{0pt}

\begin{document}

\begin{center}
	{\Large \textbf{Miraj Chamara Samarakkody}}\\[2pt]
	\href{mailto:miraj.samarakkody@gmail.com}{miraj.samarakkody@gmail.com} $\cdot$ +1 (769) 601-6290 $\cdot$ \href{https://mirajcs.github.io/}{mirajcs.github.io}
\end{center}

\section*{Research Interests}
Manifold learning and geometric methods in data science; formal proof verification in Lean 4; variational problems, geometric analysis, and mathematical physics; mathematical software development in Julia.

\section*{Education}
\textbf{Ph.D. in Mathematics}, Texas Tech University, 2021--2025\\
\textbf{B.Sc. (Hons.) Special in Mathematics}, University of Peradeniya, Sri Lanka, 2014--2018

\section*{Professional Experience}
\textbf{Assistant Professor}, Tougaloo College, Spring 2025 -- Present\\
\textbf{Research / Teaching Assistant}, Texas Tech University, 2021 -- 2024\\
\textbf{Lecturer}, University of Moratuwa, Sri Lanka, 2019 -- 2021

\section*{Selected Publications}
\begin{itemize}
	\item M. Samarakkody. ``Formalizing the Classical Isoperimetric Inequality in the Two-Dimensional Case.'' Submitted.
	\item Y. Li, E. G\"uler, M. Toda, M. Samarakkody. ``Geometric Properties of Hypersurfaces of Right Conoids in Minkowski Spacetime.'' Submitted.
	\item A. P\'ampano, M. Samarakkody, H. Tran. ``Closed $p$-Elastic Curves in Spheres of $\mathbb{L}^{3}$.'' \emph{Journal of Mathematical Analysis and Applications}, 542(2):129147, 2025.
\end{itemize}

\section*{Grants}
\begin{itemize}
	\item \textbf{PI}, NIH IDeA (P20GM103476), \$29,317 -- Statistical/ML approaches for heart disease prediction.
	\item \textbf{PI}, SCALE Program (Purdue partnership), \$73,094.77 -- Microelectronic workforce development for radiation-hardening.
	\item \textbf{Recipient}, NSF HBCU UP Implementation Project (Award No. 2510537).
\end{itemize}

\section*{Software}
\begin{itemize}
	\item \textbf{IsoperimetricInequality} -- Lean 4 formalization. \href{https://doi.org/10.5281/zenodo.20258500}{DOI: 10.5281/zenodo.20258500}
	\item \textbf{SymKit.jl} -- Julia package for symbolic computation. \href{https://doi.org/10.5281/zenodo.20294029}{DOI: 10.5281/zenodo.20294029}
	\item \textbf{ParametricCurves.jl} -- Julia package for parametric curves. \href{https://doi.org/10.5281/zenodo.20288830}{DOI: 10.5281/zenodo.20288830}
\end{itemize}

\section*{Fellowships \& Awards}
Dr. Shelby Hildebrand Graduate Math Fellowship (\$10,000, 2023--2024); AT\&T Graduate Fellowship (\$16,000, 2021--2025); University Prize for Academic Excellence, University of Peradeniya (2018).

\section*{Technical Skills}
\textbf{Languages:} Python, Julia, C \quad \textbf{Theorem Proving:} Lean 4 \quad \textbf{Math Software:} Mathematica, MATLAB \quad \textbf{Typesetting:} LaTeX, Typst

\end{document}
